% concatenated from Tract_Discv4.tex
\documentclass[12pt]{article}

\setlength{\bigskipamount}{5ex plus1.5ex minus 2ex}
\setlength{\textheight}{24cm}
\setlength{\textwidth}{16cm}
\setlength{\hoffset}{-1.3cm}
\setlength{\voffset}{-1.8cm}

\usepackage{amsmath, amssymb, amsthm, amsfonts}
\usepackage{color}
\usepackage{comment}
\usepackage{hyperref}
\usepackage{epsfig}
\usepackage{graphicx} 
%\usepackage{showlabels}


%\usepackage[utf8]{inputenc}
%\usepackage{a4wide}
\usepackage{latexsym,mathrsfs,url,bm}
%\usepackage{mathtools}
%\usepackage{lipsum}
\usepackage{xcolor}


\DeclareMathOperator{\HyperF}{F}
\DeclareMathOperator{\HyperTildeF}{\widetilde{F}}
\newcommand{\Hypergeom}[5]{{\sideset{_#1}{_#2}\HyperF\!\left(\substack{\displaystyle#3\\\displaystyle#4};#5\right)}}


\newcommand{\fritz}[1]{{\color{blue}{#1}}}
\newcommand{\josi}[1]{{\color{orange}{#1}}}
\newcommand{\hans}[1]{{\color{green}{#1}}}
\newcommand{\todo}[1]{{\color{red}{#1}}}
\newcommand{\cp}{{\rm cap}}
\newcommand{\he}{{\rm hem}}

\newtheorem{theorem}{Theorem}
\newtheorem{lemma}{Lemma}
\newtheorem{corollary}{Corollary}
\newtheorem{proposition}{Proposition}
\newtheorem{conjecture}{Conjecture}
\newtheorem{definition}{Definition}
\newtheorem{algorithm}{Algorithm}
\newtheorem{remark}{Remark}




%\title{The blessing of dimensionality of the spherical cap $L_2$ discrepancy and a balanced large-cap variant}
\title{Spherical Cap $L_2$ Discrepancy -- Blessing of Dimensionality and a Balanced Large-Cap Variant}
\author{Johann S. Brauchart\thanks{Institut für Analysis und Zahlentheorie, TU Graz, Kopernikusgasse 24/II, 8010 Graz, Austria (\url{j.brauchart@tugraz.at}). The research of J.~S.~B. was funded in whole or in part by the OeAD (Austria's Agency for Education and Internationalisation) Project No. BG 07/2025. }, Josef Dick\thanks{School of Mathematics and Statistics, The University of New South Wales Sydney, 2052 NSW, Australia (\url{josef.dick@unsw.edu.au}). The work of J.~D.\ is supported by ARC grant DP220101811.}, Friedrich Pillichshammer\thanks{Institut für Finanzmathematik und Angewandte Zahlentheorie, Johannes Kepler Universität Linz, Altenbergerstra{\ss}e 69, 4040 Linz, Austria (\url{friedrich.pillichshammer@jku.at})}}
\date{}

\begin{document}

\maketitle

\begin{abstract}
\iffalse
We show that the inverse (or information complexity) of the classical spherical cap $L_2$ discrepancy on the $d$ dimensional sphere $\mathbb{S}^d$ decreases with increasing dimension $d$, i.e., the related numerical integration problem becomes easier as the dimension increases. This behavior is unusual and may be called the ``blessing of dimensionality'' (as opposed to the ``curse of dimensionality''). We then propose a modification of the spherical cap $L_2$ discrepancy that gives more emphasis to large spherical caps and also discuss the associated reproducing kernel Hilbert space. We show that this large spherical cap $L_2$ discrepancy, and the associated numerical integration problem, do not become easier with increasing dimension. In particular, we show that $K(\bm{x}, \bm{y}) = 1 - \tfrac{1}{\sqrt{2}} \|\bm{x} - \bm{y}\|$ is a reproducing kernel for the Sobolev space $H^{(d+1)/2}$ and that the worst-case error for numerical integration in this space grows polynomially with the dimension. We also show a Stolarsky invariance principle for the large spherical cap $L_2$ discrepancy. \fi
We prove that the information complexity (i.e., the inverse) of the classical spherical cap $L_2$ discrepancy on the $d$-dimensional sphere $\mathbb{S}^d$ decreases with dimension $d$, indicating a ``blessing of dimensionality'' for the associated numerical integration problem. We then introduce a modified spherical cap $L_2$ discrepancy that emphasizes large caps (close to hemispheres). For this variant, the problem does not become easier with increasing $d$. We also establish a Stolarsky invariance principle which connects the modified spherical cap $L_2$ discrepancy to numerical integration in the Sobolev space $H^{(d+1)/2}(\mathbb{S}^d)$, represented by the reproducing kernel $K(\bm{x}, \bm{y}) = 1 - \tfrac{1}{\sqrt{2}} \|\bm{x} - \bm{y}\|$. Stolarsky's invariance principle then implies that the worst-case integration error in this space grows polynomially with $d$. %We also establish a Stolarsky invariance principle for the large spherical cap $L_2$ discrepancy, \josi{which connects the large spherical cap $L_2$ discrepancy to numerical integration in the Sobolev space $H^{(d+1)/2}$ based on the kernel $K(\bm{x}, \bm{y}) = 1 - \tfrac{1}{\sqrt{2}} \|\bm{x} - \bm{y}\|$.}.
\end{abstract}

\centerline{\begin{minipage}[hc]{130mm}{
{\em Keywords:} Spherical cap discrepancy, point distribution on the sphere, information-based-complexity, quasi-Monte Carlo\\
{\em MSC 2010:} 11K38, 65D30, 65C05, 65Y20}
\end{minipage}}


\section{Introduction}

Discrepancy on the sphere quantifies how well a finite point set reproduces the uniform surface measure when tested against indicator functions of spherical caps. Among the various notions of spherical discrepancy, the \emph{spherical cap $L_2$ discrepancy} occupies a central position: it is rotationally invariant, admits a reproducing-kernel representation, and is tied -- via Stolarsky's invariance principle -- to geometric energies on the sphere \cite{BD13a}. In numerical integration on $\mathbb{S}^d$, this discrepancy governs worst-case errors for an associated reproducing kernel Hilbert space (RKHS) of functions, thus providing an information-complexity lens on quadrature over $\mathbb{S}^d$.

This paper shows that, for the \emph{classical} spherical cap $L_2$ discrepancy on $\mathbb{S}^d$, the inverse discrepancy (equivalently, the information complexity) \emph{improves} as the ambient dimension $d$ grows: fewer points suffice to reach the same target accuracy -- the \emph{blessing of dimensionality}. This is the mirror image of the familiar ``curse of dimensionality'' in high-dimensional numerical analysis. We also show that this improvement is not a structural necessity of spherical cap-based testing. By restricting spherical caps to caps that are close to hemispheres (where ``close'' depends on the dimension $d$), we obtain a spherical cap-based discrepancy whose information complexity does \emph{not} become easier with increasing dimension $d$. We obtain a Stolarsky invariance principle for this discrepancy measure and also connect it to numerical integration on the sphere $\mathbb{S}^d$. 

\paragraph{Setting and notation.}
Let $\mathbb{S}^d:=\{ \bm{x} \in\mathbb{R}^{d+1}:\| \bm{x} \|_2=1\}$ be equipped with the normalized surface area measure $\sigma_d$, so $\sigma_d(\mathbb{S}^d)=1$, where $\| \bm{\cdot} \|_2 := \sqrt{\langle \bm{\cdot} , \bm{\cdot} \rangle}$ is the norm induced by the usual inner product in $\mathbb{R}^{d+1}$. 
%
For $\bm{x} \in\mathbb{S}^d$ and $t\in[-1,1]$ let 
\[
C(\bm{x}; t) \;:=\;\{\bm{y} \in\mathbb{S}^d:\langle \bm{x}, \bm{y} \rangle\ge t\}
\]
denote the spherical cap with center $\bm{x}$ and height $t$.

We study distribution properties of $N$-point sets $P_{N,d} :=\{\bm{x}_1,\ldots,\bm{x}_N\}$ on the sphere~$\mathbb{S}^d$. Throughout, the notation $P_{N,d}$ indicates that we speak about an $N$-element set in $\mathbb{S}^d$. A prominent quantitative measure for the uniformity of such a point set is the spherical cap $L_2$ discrepancy. In the following, let $\bm{1}_{C(\bm{x}; t)}$ denote the indicator function of a spherical cap $C(\bm{x}; t)$.



\begin{definition}\label{def1}\rm
The spherical cap $L_2$ discrepancy of an $N$-point set $P_{N,d} = \{ \bm{x}_1, \ldots, \bm{x}_N\} \subseteq \mathbb{S}^d$ is given by
\begin{equation*}
L_2^{\cp}(P_{N,d}):= \left( \int_{-1}^1 \int_{\mathbb{S}^d} \left( \frac{1}{N} \sum_{n=1}^N \bm{1}_{C(\bm{x}; t)}(\bm{x}_n) - \sigma_d(C(\bm{x}; t)) \right)^2 \,\mathrm{d} \sigma_d(\bm{x}) \,\mathrm{d} t \right)^{1/2}.
\end{equation*}
\end{definition}
As reference value we define the initial spherical cap $L_2$ discrepancy by
\begin{equation*}
L_2^{\cp}(\emptyset) := \left(\int_{-1}^1 \int_{\mathbb{S}^d} \left( \sigma_d(C(\bm{x}; t)) \right)^2 \,\mathrm{d} \sigma_d(\bm{x}) \,\mathrm{d} t\right)^{1/2}.
\end{equation*}

In addition to its geometrical interpretation, the spherical cap $L_2$ discrepancy is also intimately related to the worst-case error of quasi-Monte Carlo rules for the numerical integration of functions over $\mathbb{S}^d$ with respect to $\sigma_d$, see \cite{BD13a,BSSW14,DP14a}.



\paragraph{Information complexity.} In this paper we are interested in information-based complexity, where one typically examines the following inverse problem: \emph{determine how much information is required to reduce the initial discrepancy to an $\varepsilon$-fraction of its original value for $\varepsilon \in (0,1)$.} The corresponding measure is the information complexity, also known in the discrepancy setting as the inverse discrepancy.

\begin{definition}\rm
For $\varepsilon \in (0,1)$ and $d \in \mathbb{N}$ the inverse of the spherical cap $L_2$ discrepancy is defined as 
\begin{equation*}
N^{\cp}(\varepsilon, d) := \inf \{ N \in \mathbb{N}: \exists P_{N,d} \subseteq \mathbb{S}^d \ \text{such that}\ L_2^{\cp}(P_{N,d}) \le \varepsilon L_2^{\cp}(\emptyset) \}.
\end{equation*}
\end{definition}
Equivalently, $N^{\cp}(\varepsilon,d)$ is the minimal number of function evaluations needed by an equal-weight quadrature rule to guarantee worst-case error at most $\varepsilon$ in the unit ball of the corresponding function class. We stress that our upper and lower bounds are dimension-explicit, enabling a direct comparison of $N^{\cp}(\varepsilon,d)$ across $d$.




\paragraph{Main contributions.} The main contributions of the present work are as follows:
\begin{itemize}
  \item \textbf{Blessing of dimensionality for the classical spherical cap $L_2$ discrepancy.} We prove that, for fixed $\varepsilon\in(0,1)$, the information complexity $N^{\cp}(\varepsilon,d)$ for the classical spherical cap $L_{2}$ discrepancy is \emph{decreasing} in $d$. Thus, the associated worst-case numerical integration problem becomes easier as $d$ increases. One may call this unusual behavior the \emph{blessing of dimensionality}.

  Although it is positive from a computational effort point of view if a problem becomes easier to solve with increasing dimension, it does beg the question whether the definition of the problem in high dimensions is appropriate in the first place. We address this in the second part.
 
  \item \textbf{A large spherical cap discrepancy with stabilized difficulty.} We introduce a natural modification of the spherical cap-based $L_2$ discrepancy and its associated reproducing kernel such that the corresponding information complexity does \emph{not} exhibit a blessing of dimensionality. In this setting, $N^{\cp}(\varepsilon,d)$ does not decrease with dimension $d$. 
\end{itemize}

Besides spherical cap $L_2$ discrepancy, there are numerous other figures-of-merit for the quality of point distributions on the sphere. For at least one of these criteria, a type of diaphony, its dependence on dimension has been investigated. See, in particular, \cite[Section~3.4]{ChS13}.


\paragraph{Why the blessing of dimensionality arises.}
At a high level, our analysis isolates the following phenomenon. As $d$ grows, spherical caps undergo a concentration-of-measure effect: for a random cap anchored at an arbitrary center point $\bm{x}$, its surface area measure concentrates around the equator (relative to $\bm{x}$) at a rate controlled by $d$. More precisely, for fixed $t \in (0,1)$ and any $\bm{x} \in \mathbb{S}^d$ we have (see Section~\ref{sec_dist_int})

\begin{equation}\label{asym:Cxt}
\sigma_d(C(\bm{x}; t)) \sim \frac{1}{\sqrt{2\pi d}} \, \frac{(1-t^2)^{d/2}}{t} = \frac{1}{\sqrt{2\pi d}} \, \frac{o\big( \mathrm{e}^{- d t^2 / 2} \big)}{t} \qquad \text{as $d \to \infty$.}
\end{equation}
This means that the measure of spherical caps decreases exponentially fast with increasing dimension when $t \in (0,1)$ remains fixed (or, more generally, $d t^2$ tends to infinity sufficiently fast). However, for $t = 0$ (which corresponds to a hemisphere), the measure is $1/2$ and for $t$ close enough to $0$ (depending on $d$) the measure is also close to $1/2$. Since the $L_2$ spherical cap discrepancy takes the average of the square of the local discrepancy function over $t$ from $-1$ to $1$, for increasing dimension $d$, more and more weight is given to spherical caps which are very small in surface measure. In the limit, as $d \to \infty$, the discrepancy measure of a point set becomes trivial (only spherical caps close to the hemisphere ($t=0$) matter in this case). In order to avoid that problem for large dimensions, we propose to restrict $t$ to the interval $[-\alpha_d, \alpha_d]$ for an appropriately chosen $\alpha_d \in [0,1]$ depending on the dimension such that spherical caps which are ``too small'' in volume get ignored. 

\paragraph{The large spherical cap $L_2$ discrepancy.}

In the following, we introduce a parameter in the definition of the spherical cap $L_2$ discrepancy which will allow us to adjust the behavior of the inverse of the $L_2$ spherical cap discrepancy as the dimension increases (see the discussion in Section~\ref{sec:tract}).

%We follow \cite{BD13a}. 
\newpage
\begin{definition}\rm
For $d \in \mathbb{N}$ let $\alpha_d \in [0, 1]$. The $\alpha_d$-large spherical cap $L_2$ discrepancy of an $N$-point set $P_{N,d} =\{ \bm{x}_1, \ldots, \bm{x}_N\} \subseteq \mathbb{S}^d$ is given by
\begin{equation*}
L_2^{\cp,a_d}(P_{N,d}) :=\left(\int_{-1}^1 \int_{\mathbb{S}^d} \left|\frac{1}{N} \sum_{n=1}^N \bm{1}_{C(\bm{x}; \alpha_d \, t)}(\bm{x}_n) -\sigma_d(C(\bm{x}; \alpha_d \, t))\right|^2 \,\mathrm{d} \sigma_d(\bm{x}) \,\mathrm{d} t\right)^{1/2}.
\end{equation*}
\end{definition}

Choosing $\alpha_d=1$ yields Definition~\ref{def1}. For $-1 \le t \le 1$ we have $-\alpha_d \le t\alpha_d \le \alpha_d$. 
%
For small $\alpha_d$ and in particular when $\alpha_d \to 0$ as $d \to \infty$ the large spherical cap $L_2$ discrepancy uses spherical caps that are close to a hemisphere.
%
We discuss the extreme case $\alpha_d = 0$ separately in Section~\ref{sec_hemisphere}.
%
An interesting special case is where $\alpha_d = \sqrt{2} \, C_d$, in which case we obtain polynomial tractability of the $\alpha_d$-large spherical cap $L_2$ discrepancy and a reasonable normalization of the corresponding integration problem, see Section~\ref{sec:tract}. %, in particular Eq.~\eqref{eq_alphadchoice}.


We note that, in principle, one could introduce a function $s:[-1,1] \to [-1,1]$ and consider caps $C(\bm{x}; s(t))$ to avoid the blessing of dimensionality. In our case, we only consider the special case $s(t) = \alpha_d \, t$. The reason for our particular choice is that with $s(t) = \alpha_d \, t$ we can still obtain a Stolarsky invariance principle and can connect the $\alpha_d$-large spherical cap $L_2$ discrepancy to numerical integration on the sphere $\mathbb{S}^d$. Though other choices of functions $s$ may turn out to be useful in the future, a version of Stolarsky's invariance principle and its connection to numerical integration on the sphere for general functions $s$ does not appear to be obvious.

We discuss appropriate choices for $\alpha_d$ in the ``Discussion'' section of Section~\ref{sec:tract}. There we identify choices of $\alpha_d$ such that the inverse of the large spherical cap $L_2$ discrepancy problem does not get easier with increasing dimension $d$.

\paragraph{Stolarsky's invariance principle and numerical integration.}

We return to the classical setting. 
%
In \cite{BD13a} the first two authors show that for the Sobolev space $H^{(d+1)/2}$ over the sphere $\mathbb{S}^d$ provided with the reproducing kernel
\begin{equation*}
K_{1,d}(\bm{x}, \bm{y}) := 1 - C_d \|\bm{x} - \bm{y}\|, \qquad \bm{x}, \bm{y} \in \mathbb{S}^d, 
\end{equation*}
where (see \cite[Eq. (2.4)]{BD13a})
\begin{equation}\label{def:Cd}
C_d := \frac{1}{d} \frac{\Gamma(\tfrac{d+1}{2})}{\sqrt{\pi} \ \Gamma(\tfrac{d}{2})} \sim \frac{1}{\sqrt{2 \pi d}} \quad \mbox{as } d \to \infty
\end{equation}
expressed in terms of the Gamma function $\Gamma(s) := \int_0^{\infty} t^{s-1} {\rm e}^{-t} \, \mathrm{d} t$, $\Re s > 0$, the spherical cap $L_2$ discrepancy of an $N$-point set $P_{N,d}$ and the worst-case numerical integration error of an equal weight numerical integration rule taking the average of function values at the nodes in $P_{N,d}$ for functions in $H^{(d+1)/2}$ have the same value. 
%
Furthermore, the squared spherical cap $L_2$ discrepancy can be computed either using Stolarsky's invariance principle 
\begin{equation}
\label{eq:Stolarsky.principle}
(L_2^{\cp}(P_{N,d}))^2 = C_d  \int_{\mathbb{S}^d} \int_{\mathbb{S}^d} \|\bm{x} - \bm{y}\| \,\mathrm{d} \sigma_d(\bm{x}) \,\mathrm{d}\sigma_d(\bm{y}) - C_d \frac{1}{N^2} \sum_{m,n=1}^N \|\bm{x}_m -\bm{x}_n\|    
\end{equation}
or the well-known reproducing kernel Hilbert space representation of the squared  worst case error
\begin{align} \label{L2_kernel}
\begin{split}
(L_2^{\cp}(P_{N,d}))^2
&=
\int_{\mathbb{S}^d} \int_{\mathbb{S}^d} K_{1,d}(\bm{x}, \bm{y}) \,\mathrm{d} \sigma_d(\bm{x}) \,\mathrm{d}\sigma_d(\bm{y}) 
-
\frac{2}{N} \sum_{n=1}^N \int_{\mathbb{S}^d} K_{1,d}(\bm{x}, \bm{x}_n) \,\mathrm{d} \sigma_d(\bm{x}) 
\\
&\phantom{=}+
\frac{1}{N^2} \sum_{m,n=1}^N K_{1,d}(\bm{x}_m, \bm{x}_n).
\end{split} 
\\
&=
\frac{1}{N^2} \sum_{m,n=1}^N K_{1,d}(\bm{x}_m, \bm{x}_n) - \int_{\mathbb{S}^d} \int_{\mathbb{S}^d} K_{1,d}(\bm{x}, \bm{y}) \,\mathrm{d} \sigma_d(\bm{x}) \,\mathrm{d}\sigma_d(\bm{y}). \notag
\end{align}
%
Formula \eqref{L2_kernel} provides a way to compute the initial spherical cap $L_2$ discrepancy by observing that for an empty point set the last two sums in \eqref{L2_kernel} are not present and thus
\begin{equation}\label{init:err:fo}
(L_2^{\cp}(\emptyset))^2 = \int_{\mathbb{S}^d} \int_{\mathbb{S}^d} K_{1,d}(\bm{x}, \bm{y}) \,\mathrm{d} \sigma_d(\bm{x}) \,\mathrm{d} \sigma_d(\bm{y}) = 1 - C_d I_d,
\end{equation}
where $I_d$ denotes the distance integral given as
\begin{equation}\label{def:Id}
I_d :=  \int_{\mathbb{S}^d} \int_{\mathbb{S}^d} \|\bm{x} - \bm{y}\| \,\mathrm{d} \sigma_d(\bm{x}) \,\mathrm{d} \sigma_d(\bm{y}).
\end{equation}

Before we continue with our explanations, we collect some estimates of the quantities $I_d$ and $C_d$ in the following lemma. The proof is deferred to Appendix~\ref{sec_dist_int}. 
\begin{lemma}
\label{lem:C.d.I.d.inequalities}
The constants $I_d$ of \eqref{def:Id} and $C_d$ of \eqref{def:Cd} satisfy for all dimensions $d \in \mathbb{N}$ the following relations:
\begin{enumerate}
\item $\displaystyle \sqrt{2} \, \sqrt{1 - \frac{1}{d}} < I_d < \sqrt{2}$; furthermore, $\displaystyle I_d = \sqrt{2} \left( 1 - \tfrac{1}{8d} + \mathcal{O}\big( \tfrac{1}{d^2} \big) \right)$ as $d \to \infty$;
\item $\displaystyle \frac{1}{\sqrt{2 \pi \, d}} \, \sqrt{1 - \frac{1}{d}} < C_d < \frac{1}{\sqrt{2 \pi \, d}}$;
\item $\displaystyle \frac{1}{\sqrt{\pi \, d}} \left( 1 - \frac{1}{d} \right) < C_d \, I_d < \frac{1}{\sqrt{\pi \, d}}$;
\item $\displaystyle  \frac{1 - \frac{1}{d}}{\sqrt{\pi \, d} - \left( 1 - \frac{1}{d} \right)} < \frac{C_d \, I_d}{1 - C_d \, I_d} < \frac{1}{\sqrt{\pi \, d} - 1}$.
\end{enumerate}
\end{lemma}

From Lemma~\ref{lem:C.d.I.d.inequalities} one can easily deduce leading term asymptotics for the quantities $I_d$, $C_d$, $C_d \, I_d$, and $C_d \, I_d / ( 1 - C_d \, I_d )$ as $d \to \infty$.

Now we will continue with our explanations about the spherical cap $L_2$ discrepancy. Applying Lemma~\ref{lem:C.d.I.d.inequalities} to the definition of the reproducing kernel we obtain 
\begin{equation}\label{kernel_bounds}
1 - 2 C_d \le K_{1,d}(\bm{x}, \bm{y}) \le 1 
\end{equation}
and $K_{1,d}(\bm{x}, \bm{y}) \to 1$ as $d \to \infty$.

The distance integral has the explicit expression (see, e.g., \cite{BSSW14})
\begin{equation}
\label{eq:I.d.value}
I_d = 2^d \, \frac{\Gamma( \tfrac{d + 1}{2} ) \, \Gamma( \tfrac{d + 1}{2} )}{\sqrt{\pi} \, \Gamma( d + \tfrac{1}{2} )}
%=
%\sqrt{2} \, \frac{\Gamma( (d/2) + (1/2) )}{\Gamma( (d/2) + (1/4) )} \, \frac{\Gamma( (d/2) + (1/2) )}{\Gamma( (d/2) + (3/4) )} 
\sim
\sqrt{2} 
\qquad
\text{as $d \to \infty$,}
\end{equation}
where the asymptotic formula is obtained from Lemma~\ref{lem:C.d.I.d.inequalities}. Since, again by Lemma~\ref{lem:C.d.I.d.inequalities}, $C_d \, I_d \sim 1/\sqrt{\pi d}$ as $d \to \infty$, we observe
\begin{equation*}
L_2^{\cp}(\emptyset) = \sqrt{1 - C_d \, I_d} \sim \sqrt{ 1 - \frac{1}{\sqrt{\pi d}}} \sim 1 - \frac{1}{\sqrt{4\pi d}} \qquad \text{as $d \to \infty$.}
\end{equation*}

The blessing of dimensionality can also be observed from the reproducing kernel $K_{1,d}$. As noted above, as $d \to \infty$, we have $C_d \to 0$ and $K_{1,d} \to 1$. This means that the corresponding reproducing kernel Hilbert space converges to a space of only constant functions (which can be integrated exactly using just one point). Again, this illustrates that the integration problem becomes easier as the dimension $d$ increases.


\paragraph{The reproducing kernel corresponding to the large spherical cap $L_2$ discrepancy.} As will be shown in Section~\ref{sec:tract} (in particular \eqref{Kalpha}), the reproducing kernel associated with the large spherical cap $L_2$ discrepancy is given by
\begin{equation*}
K_{\alpha_d, d}(\bm{x}, \bm{y}) = 1 - \frac{C_d}{\alpha_d} \|\bm{x} - \bm{y}\|, \qquad \bm{x},\bm{y} \in \mathbb{S}^d.
\end{equation*}
Arguably, the most interesting special case is
\begin{equation*}
\alpha_d = \sqrt{2} \, C_d,
\end{equation*}
see again Section~\ref{sec:tract}, in particular Eq.~\eqref{eq_alphadchoice}. In this case, the reproducing kernel 
\begin{equation}\label{Ksqrt2Cd}
K_{\sqrt{2} C_d,d}(\bm{x}, \bm{y}) = 1 - \frac{1}{\sqrt{2}} \|\bm{x} - \bm{y}\|, \qquad \bm{x},\bm{y} \in \mathbb{S}^d
\end{equation}
does not converge to $1$ as $d \to \infty$ and hence the corresponding reproducing kernel Hilbert space does not converge to a space of only constant functions and the associated numerical integration problem becomes not easier with increasing dimension. In fact, we will show in \eqref{Nsqrt2Cd} that for this choice the inverse discrepancy grows with the dimension with order at least $\sqrt{d}$ and at most $d$.


\paragraph{Comparison with the discrepancy on the unit cube.}

The canonical example is the $L_2$ discrepancy of $N$-point sets in the $d$-dimensional unit cube $[0,1)^d$ with respect to the uniform (Lebesgue) measure $\lambda$ given by
\begin{equation}\label{def:L2discCube}
L_2(P_{N,d})=\left(\int_{[0,1]^d} \left| \frac{1}{N} \sum_{n=1}^N \bm{1}_{[\bm{0}, \bm{t})}(\bm{x}_n) - \lambda([\bm{0}, \bm{t})) \right|^2 {\rm d} \bm{t}\right)^{1/2},
\end{equation}
where $P_{N,d}=\{\bm{x}_1,\ldots,\bm{x}_N\} \subseteq [0,1)^d$, $[\bm{0},\bm{t})$ denotes the axis-parallel box anchored at the origin with opposite corner $\bm{t}\in[0,1]^d$, and $\bm{1}_{[\bm{0},\bm{t})}$ is the indicator function of the box $[\bm{0},\bm{t})$. Similarly as for spheres, the $L_2$ discrepancy on the unit cube coincides with the worst-case error of quasi-Monte Carlo (QMC) rules for numerical integration. For background, see \cite[Chapter~9]{NW10}.

For the $L_2$ discrepancy on the unit cube \eqref{def:L2discCube} it is known that the number of points required to reduce the initial discrepancy by a factor of $\varepsilon$ increases exponentially with the dimension -- a behavior termed ``curse of dimensionality'', see, e.g., \cite{NP25b,NW10,Wo99}. 

Various other notions of how $N(\varepsilon,d)$ depends on $d$ have been discussed in the literature. If $N(\varepsilon, d)$ grows at most polynomially in $d$ and $\varepsilon^{-1}$, the problem is \emph{polynomially tractable}; if the growth is at most polynomial in $\varepsilon^{-1}$ uniformly in $d$, one speaks of \emph{strong polynomial tractability}. Recent years have seen many discrepancy notions classified according to the growth of their inverse; see, e.g., \cite{NP25a,NW08,NW10}.




\paragraph{Consequences and examples.}
The blessing-of-dimensionality result has concrete implications. In applications that naturally live on high-dimensional spheres -- such as directional statistics for very high-dimensional data, numerical integration of isotropic kernels in machine learning on spheres, or Monte Carlo rendering on spherical domains -- the classical spherical cap $L_2$ discrepancy provides optimistic sample-complexity predictions as $d$ grows. When such optimism is undesirable, our modified discrepancy supplies a dimension-robust alternative with essentially the same algorithmic primitives.






\section{The blessing of dimensionality of the spherical cap $L_2$~discrepancy}\label{sec:scd1}

The following result shows the blessing of dimensionality for the classical spherical cap $L_2$ discrepancy. 

\begin{theorem}
\label{thm:1}
For all dimensions $d \in \mathbb{N}$ holds
\begin{equation*}
N^{\cp}(\varepsilon, d) \leq \left\lceil \frac{1}{\varepsilon^2} \, \frac{1}{\sqrt{\pi  d} - 1} \right \rceil.
\end{equation*}
\end{theorem}

\begin{proof}
Assume that the points in $P_{N, d} = \{\bm{x}_1, \ldots, \bm{x}_N\}$ are chosen randomly i.i.d. uniformly distributed on $\mathbb{S}^d$. Then using Stolarsky's invariance principle \eqref{eq:Stolarsky.principle} and \eqref{def:Id}
\begin{equation}\label{exp_rand}
\begin{split}
\mathbb{E} \left[ (L^{\cp}_2(P_{N,d}))^2 \right] 
&= 
C_d \left( I_d - \frac{1}{N^2} \sum_{m,n=1}^N \mathbb{E}\left[ \| \bm{x}_m - \bm{x}_n\| \right] \right) \\
&= 
C_d \left( I_d - \frac{N^2-N}{N^2} I_d \right) = \frac{C_d I_d}{N},
\end{split}
\end{equation}
where $C_d$ is defined in \eqref{def:Cd}.\footnote{We note that in \cite[Section~7]{BSSW14} the expected value for the squared worst-case error for a whole class of reproducing kernels is obtained.} From \eqref{exp_rand} it follows that for each dimension $d$ and each natural number $N$, there exists a point set $P_{N,d}$ such that
\begin{equation*}
L_2^{\cp}(P_{N,d}) \le \sqrt{\frac{C_d I_d}{N}}.
\end{equation*}
Comparing with the initial spherical cap $L_2$ discrepancy \eqref{init:err:fo}, we find that $N^{\cp}(\varepsilon, d) \le N$ for any $N \in \mathbb{N}$ which satisfies
\begin{equation*}
\sqrt{\frac{C_d I_d}{N (1-C_d I_d)}} \le \varepsilon \quad \mbox{ or, equivalently,}\quad N \ge \frac{1}{\varepsilon^2} \frac{C_d I_d}{1-C_d I_d}.
\end{equation*}
Therefore
\begin{equation*}
N^{\cp}(\varepsilon, d) \le \left\lceil \frac{1}{\varepsilon^2} \frac{C_d I_d}{1-C_d I_d} \right\rceil.
\end{equation*}
The result follows by item 4. of Lemma~\ref{lem:C.d.I.d.inequalities}.
\end{proof}

Theorem~\ref{thm:1} means that the integration problem becomes easier as the dimension $d$ increases, since $N^{\cp}(\varepsilon, d)$ decreases with increasing dimension $d$. This can be understood by considering \eqref{kernel_bounds}. As $d \rightarrow \infty$, the kernel $K_{1,d}$ approaches the constant function $1$, which means that the corresponding reproducing kernel Hilbert space approaches the space of constant functions. For instance, for the special case where $N = 1$, due to symmetry, the location of $\bm{x}_1$ is irrelevant. Then for any fixed $\bm{x}_1 \in \mathbb{S}^d$ we have
\begin{equation*}
\mathbb{E}\left[(L_2^{\cp}(P_{1,d}))^2 \right] = (L_2^{\cp}(\{\bm{x}_1\}))^2 = C_d I_d \sim \frac{1}{\sqrt{\pi d}} \quad \mbox{ as } d \to \infty.
\end{equation*}
Since a constant function can be integrated exactly using only one point, this illustrates that the space of functions converges to the space of constant functions.




\section{Tractability properties of the large spherical cap $L_2$ discrepancy}\label{sec:tract}

\paragraph{Stolarsky's invariance principle for the large spherical cap discrepancy.}
For $\alpha_d \in [0, 1]$ we also define
\begin{equation*}
K_{\alpha_d,d}(\bm{x}, \bm{y}) := \int_{-1}^1 \int_{\mathbb{S}^d} \bm{1}_{C(\bm{z}; \alpha_d \, t)}(\bm{x}) \bm{1}_{C(\bm{z}; \alpha_d \, t)}(\bm{y}) \,\mathrm{d} \sigma_d(\bm{z}) \,\mathrm{d} t.
\end{equation*}
Note that $K_{\alpha_d,d}(\bm{x}, \bm{y}) = K_{\alpha_d,d}(\bm{y}, \bm{x})$ and for any $N \in \mathbb{N}$, any $c_1,\ldots,c_N \in \mathbb{C}$ and any $\bm{x}_1,\ldots,\bm{x}_N \in \mathbb{S}^d$ we have
\begin{equation*}
\sum_{n,m=1}^N c_n \overline{c}_m K_{\alpha_d,d}(\bm{x}_n, \bm{x}_m) = \int_{-1}^1 \int_{\mathbb{S}^d} \left| \sum_{n=1}^N c_n \bm{1}_{C{(\bm{z}; \alpha_d \, t)}}(\bm{x}_n) \right|^2 \,\mathrm{d} \sigma_d(\bm{z})\,\mathrm{d} t \ge 0.
\end{equation*}
Thus, $K_{\alpha_d,d}$ is a reproducing kernel.

Using the approach from \cite[Section~2]{BD13a} we obtain for any $\alpha_d \in (C_d I_d, 1]$ that
\begin{equation}\label{Kalpha}
K_{\alpha_d,d}(\bm{x}, \bm{y}) = 1 - \frac{C_d}{\alpha_d} \left\|\bm{x} - \bm{y} \right\|.
\end{equation}

\begin{remark}\rm\label{rm_rk}
For $\alpha_d \in (C_d I_d,1]$ we have that the kernel $K_{\alpha_d,d}$ can be expanded as a series in spherical harmonics with only diagonal terms and whose coefficients are all positive. For the constant part of the kernel this follows from item 1. in Lemma~\ref{lem:C.d.I.d.inequalities} and for the remaining coefficients in the Karhunen-Leove expansion this follows from \cite{BD13a}. This implies that $K_{\alpha_d,d}$ is a reproducing kernel for the Sobolev space $H^{(d+1)/2}$, which is the space defined by the classical kernel $K_{1,d}$. Thus for any $\alpha_d \in (C_d I_d, 1]$, the numerical integration problem is based on the same function space $H^{(d+1)/2}$, albeit with a different (but equivalent) norm, which exhibits a different dependence on the dimension. This applies, in particular, to $\alpha_d = \sqrt{2} \, C_d \in (C_d I_d, 1]$, which leads to the reproducing kernel \eqref{Ksqrt2Cd}. Note that $C_d \, I_d < \sqrt{2} \, C_d \le \sqrt{2\pi} \, C_d \le 1$, using Lemma~\ref{lem:C.d.I.d.inequalities}.
\end{remark}

For any $\alpha_d \in (0,1]$, the corresponding Stolarsky invariance principle is
\begin{equation}
\label{eq:Stolarsky.gen}
\begin{split}
(L_2^{\cp,a_d}(P_{N,d}))^2 =&  \frac{1}{N^2} \sum_{n,m=1}^N K_{\alpha_d,d}(\bm{x}_n, \bm{x}_m) - \int_{\mathbb{S}^d} \int_{\mathbb{S}^d} K_{\alpha_d,d}(\bm{x}, \bm{y}) \,\mathrm{d} \sigma_d(\bm{x}) \,\mathrm{d} \sigma_d(\bm{y}).
\end{split}
\end{equation}


\begin{remark}\rm
By Stolarsky's invariance principle, point sets $\{\bm{x}_1^*, \ldots, \bm{x}_N^*\}$ which maximize $\sum_{n,m=1}^N \|\bm{x}_n - \bm{x}_m\|$, minimize the spherical cap $L_2$ discrepancy. By \eqref{eq:Stolarsky.gen}, such point sets also minimize the $\alpha_d$-large spherical cap $L_2$ discrepancy for any $\alpha_d \in (0,1]$.
\end{remark}

\paragraph{Tractability of a normalized kernel.} As reference value, we define the initial large spherical cap $L_2$ discrepancy by
\begin{equation*}
L_{2,\alpha_d}(\emptyset) = \left(\int_{\mathbb{S}^d} \int_{\mathbb{S}^d} K_{\alpha_d, d}(\bm{x}, \bm{y}) \,\mathrm{d} \sigma_d(\bm{x}) \,\mathrm{d} \sigma_d(\bm{y})\right)^{1/2} = \left(1 - \frac{C_d I_d}{\alpha_d}\right)^{1/2}.
\end{equation*}

\begin{definition}\rm
For $\varepsilon > 0$ and $d \in \mathbb{N}$ the inverse of the large spherical cap $L_2$ discrepancy is defined as 
\begin{equation*}
N^{\cp,\alpha_d}(\varepsilon, d) := \inf \{ N \in \mathbb{N}: \exists P_{N,d} \subseteq \mathbb{S}^d \ \text{such that}\ L_2^{\cp,\alpha_d}(P_{N,d}) \le \varepsilon L_2^{\cp,\alpha_d}(\emptyset) \}.
\end{equation*}
\end{definition}

\begin{theorem}\label{thm1}
For all $\varepsilon \in (0,1)$ and all $d \in \mathbb{N}$, we have
\begin{equation}\label{Ned_alpha}
\frac{1}{8} \left( \frac{1}{\sqrt{d}} \, \frac{1}{\varepsilon^2} \, \frac{C_d \, I_d}{\alpha_d - C_d \, I_d} \right)^{\tfrac{d}{d+1}} 
\leq 
N^{\cp,\alpha_d}(\varepsilon, d) 
\leq \left\lceil\frac{1}{\varepsilon^2} \, \frac{C_d \, I_d}{\alpha_d - C_d \, I_d}\right\rceil.
\end{equation}
\end{theorem}

\begin{proof}
The average large spherical cap $L_2$ discrepancy of a uniformly distributed point set on the sphere satisfies
\begin{align*}
\mathbb{E}\left[(L_2^{\cp,\alpha_d}(P_{N,d}))^2\right] & = \frac{1}{N} \left( \int_{\mathbb{S}^d} K_{\alpha_d, d}(\bm{x}, \bm{x}) \,\mathrm{d} \sigma_d(\bm{x}) - \int_{\mathbb{S}^d} \int_{\mathbb{S}^d} K_{\alpha_d, d}(\bm{x}, \bm{y}) \,\mathrm{d} \sigma_d(\bm{x}) \, \mathrm{d} \sigma_d(\bm{y}) \right) \\ & = \frac{1}{N} \left( 1 - 1 + \frac{C_d I_d}{\alpha_d}  \right) = \frac{C_d I_d}{\alpha_d N}.
\end{align*}
As in the proof of Theorem~\ref{thm:1}, this shows that
\begin{equation*}
N^{\cp,\alpha_d}(\varepsilon, d) \le \left\lceil \frac{1}{\varepsilon^2} \, \frac{C_d I_d}{\alpha_d - C_d I_d }\right\rceil.
\end{equation*}
Hence, the upper bound in \eqref{Ned_alpha} is shown.

The lower bound in \eqref{Ned_alpha} follows from the general lower bound (shown in Appendix~\ref{app:lower.bound})
\begin{equation}\label{L2_lower_bound}
(L_2^{\cp,\alpha_d}(P_{N,d}))^2  \geq 
\frac{C_d I_d}{\alpha_d}  \, \frac{F(x_d^*)}{2 \sqrt{2\pi}} \, \frac{1}{N^{(d+1)/d}},
\end{equation}
where (cf. Eq.~\eqref{eq:F.x.star})
\begin{equation*}
F(x_d^*) = \frac{d}{(d+1)^{\tfrac{d+1}{d}}} \left(\frac{\Gamma(\frac{1}{2})}{\Gamma(\frac{d+1}{2})}\right)^{\tfrac{1}{d}} 
\end{equation*}
for any $N$-point set $P_{N,d}$ on $\mathbb{S}^d$. This implies that 
\begin{equation*}
N^{\cp,\alpha_d}(\varepsilon, d) 
\geq 
\left( \frac{1}{\varepsilon^2} \, \frac{C_d \, I_d}{\alpha_d - C_d \, I_d} \right)^{\tfrac{d}{d+1}} \left( \frac{F(x_d^*)}{2 \sqrt{2 \pi}} \right)^{\tfrac{d}{d+1}}.
\end{equation*}
We use \cite[Eq.~5.6.1]{NIST}, i.e. $\Gamma( x ) < \sqrt{2 \pi } \, x^{x-1/2} {\rm e}^{-x} \, {\rm e}^{1/(12x)}$ for $x \in \mathbb{R}^+$, and the fact that $\Gamma(\tfrac{1}{2})=\sqrt{\pi}$ to get
\begin{align*}
\left( \frac{F(x_d^*)}{2 \sqrt{2\pi}} \right)^{\tfrac{d}{d+1}}
&=
\frac{d^{\tfrac{d}{d+1}}}{d+1} \left( \frac{1}{2 \sqrt{2 \pi}} \right)^{\tfrac{d}{d+1}} \left(\frac{\Gamma(\frac{1}{2})}{\Gamma(\frac{d+1}{2})}\right)^{\tfrac{1}{d+1}}
\\
&>
\frac{1}{2 \sqrt{2 \pi}} \, \frac{d}{d+1} \left( \frac{2 \, \sqrt{2\pi}}{d} \right)^{\tfrac{1}{d+1}} \left(\frac{\sqrt{\pi}}{\sqrt{2\pi} \left( \frac{d+1}{2} \right)^{\tfrac{d}{2}} {\rm e}^{-\tfrac{d+1}{2}} \, {\rm e}^{\tfrac{1}{6(d+1)}}}\right)^{\tfrac{1}{d+1}}
\\
%&=
%\frac{\sqrt{{\rm e}}}{2 \sqrt{2 \pi}} \left( \frac{2 \, \sqrt{\pi}}{d} \right)^{\tfrac{1}{d+1}} \left( \frac{\sqrt{2}}{\sqrt{d+1}} \right)^{\tfrac{d}{d+1}} \frac{{\rm e}^{-\tfrac{1}{6(d+1)^2}}}{1 + \frac{1}{d}}
%\\
&= \frac{\sqrt{{\rm e}}}{2 \, \sqrt{\pi}} \left( \frac{1}{\sqrt{d}} \right)^{\tfrac{d}{d+1}}  f(d) \, h(d),
\end{align*}
where 
\begin{equation*}
f(d) := \left( \frac{\sqrt{2 \pi}}{d} \sqrt{1 + \frac{1}{d}}\right)^{\tfrac{1}{d+1}} \quad \mbox{ and } \quad h( d ) := \frac{{\rm e}^{-\tfrac{1}{6(d+1)^2}}}{\left( 1 + \frac{1}{d} \right)^{\tfrac{3}{2}}}.
\end{equation*}
%
We observe that the function $f$ has in $\mathbb{R}^+$ a unique minimum somewhere in $(8,9)$ such that $f(d) \geq f( 9 )$ for all $d \in \mathbb{N}$. Furthermore, the function $h$ is strictly monotonically increasing (as can be seen from its positive first derivative) in $\mathbb{R}^+$ such that $h(d) \geq h(1)$ for all $d \in \mathbb{N}$. Thus 
\begin{equation*}
N^{\cp,\alpha_d}(\varepsilon, d) 
\geq 
c \left( \frac{1}{\sqrt{d}} \, \frac{1}{\varepsilon^2} \, \frac{C_d \, I_d}{\alpha_d - C_d \, I_d} \right)^{\tfrac{d}{d+1}},
\end{equation*}
where 
\begin{equation*}
c := \frac{\sqrt{{\rm e}}}{2 \, \sqrt{\pi}} \, f(9) \, h(1) 
= 
\frac{5^{1/20} {\rm e}^{11/24}}{4 \cdot 2^{2/5} 3^{3/10} \pi^{9/20}}
=
0.1395\ldots 
>
\frac{1}{8}.
\end{equation*}

%
This finishes the proof of Theorem~\ref{thm1}.
\end{proof}

\paragraph{Discussion.} First we note that since $C_d I_d \le 1/\sqrt{\pi d}$, by  Lemma~\ref{lem:C.d.I.d.inequalities}, we have the blessing of dimensionality whenever $\alpha_d - C_d I_d \ge c > 0$ with an absolute constant $c$.

 


Now let $g:\mathbb{N} \to (0, 1)$ be a given function (we discuss particular choices below). By choosing $\alpha_d$ such that $1- I_d C_d / \alpha_d = g(d)$, i.e.
\begin{equation}\label{def_alphad}
\alpha_d = \frac{I_d C_d}{1- g(d)},
\end{equation}
the initial large spherical cap $L_2$ discrepancy is $L_2^{\cp,\alpha_d}(\emptyset) = \sqrt{g(d)}$. In order for the corresponding spaces to be equivalent to the Sobolev space $H^{(d+1)/2}$, we need to ensure that $1-\frac{C_d}{\alpha_d} I_d > 0$. This means that we need $\alpha_d > C_d I_d$, which is readily satisfied for $\alpha_d = I_d C_d / (1-g(d))$.

%The average of the square discrepancy over $N$ random points is
%\begin{align*}
%\mathbb{E}(L_{2,\alpha_d}^2(P_{N,d}) & = \frac{1}{N} \left( 1 - 1 + \frac{C_d I_d}{\alpha_d} \right) = \frac{1- g(d)}{N}.
%\end{align*}

From \eqref{def_alphad} we obtain $C_d I_d/(\alpha_d- C_d I_d)=(1-g(d))/g(d)$. Substituting into Theorem~\ref{thm1} we obtain 
\begin{equation*}
\frac{1}{8} \left( \frac{1}{\sqrt{d}} \frac{1}{\varepsilon^2} \frac{1-g(d)}{g(d)} \right)^{\tfrac{d}{d+1}}  \le N^{\cp,\alpha_d}(\varepsilon, d) \le \left\lceil \frac{1}{\varepsilon^2} \frac{1-g(d)}{g(d)}\right\rceil.
\end{equation*}


For different choices of $g$ we observe rather different behaviors:
\begin{itemize}
\item If $g(d) \to 1$ as $d \to \infty$, then we get the blessing of dimensionality. This case comprises the classical spherical cap $L_2$ discrepancy from Section~\ref{sec:scd1}. 


\item If $g(d) \to 0^+$ as $d \to \infty$, then we get from \eqref{def_alphad} that $\alpha_d \sim C_d I_d$ as $d\to \infty$. However, the dimension dependence can vary wildly depending on the decay rate of $g(d)$. In other words, $N^{\cp,\alpha_d}(\varepsilon, d)$ is very sensitive to small changes in $\alpha_d$. This phenomenon may be explained by \eqref{asym:Cxt}, which shows that the measure of spherical caps changes exponentially fast with the dimension, and so the measure of spherical caps $C(\bm{x};t)$ is very sensitive to $t$ for large dimensions.
\begin{itemize}
\item An important special case is $g(d)=1-I_d/\sqrt{2}$, which means $\alpha_d = \sqrt{2} C_d \in (I_d C_d, 1]$ by Lemma~\ref{lem:C.d.I.d.inequalities}. Then $g(d)\rightarrow 0^+$ as $d \to \infty$. According to Lemma~\ref{lem:C.d.I.d.inequalities} we have  
\begin{equation*}
\alpha_d - C_d I_d = \sqrt{2} C_d \left(1 - \frac{I_d}{\sqrt{2}}\right) \sim  \frac{\sqrt{2} C_d}{8d} \qquad \mbox{ as } d \to \infty
\end{equation*}
such that $$\frac{C_d I_d}{\alpha_d - C_d I_d} \sim \frac{\sqrt{2} C_d}{\alpha_d - C_d I_d} \sim 8 d  \qquad \mbox{ as } d \to \infty.$$
Hence, Theorem~\ref{thm1} gives
\begin{equation}\label{Nsqrt2Cd}
\left( \frac{ \sqrt{d} }{\varepsilon^2}  \right)^{\tfrac{d}{d+1}} \lesssim N^{\cp,\sqrt{2} C_d}(\varepsilon, d) \lesssim \frac{ d }{\varepsilon^2}.
\end{equation}
(The sign $\lesssim$ means that we suppress absolute positive constants in the notation.) In particular, we do not have the blessing of dimensionality in this case, but polynomial tractability. 

We recall that by choosing $\alpha_d = \sqrt{2} C_d$, we obtain 
\begin{equation}\label{eq_alphadchoice}
K_{\sqrt{2} C_d,d}(\bm{x}, \bm{y}) = 1 - \frac{1}{\sqrt{2}} \|\bm{x} - \bm{y}\|. 
\end{equation}
In this case, the constant $1/\sqrt{2}$ does not depend on the dimension, and so the kernel function $K_{\sqrt{2} C_d, d}$ does not tend to 1 as $d \to \infty$ (compare with~\eqref{kernel_bounds}). This kernel is also a reproducing kernel for the Sobolev space $H^{(d+1)/2}$ by Remark~\ref{rm_rk}.



\item If $g(d)$ decays polynomially fast to 0, say $g(d) \asymp d^{-\delta}$ as $d\to \infty$, then we do not obtain the blessing of dimensionality, but we obtain at least polynomial tractability. If $\delta >1/2$ we cannot obtain strong polynomial tractability. 
\item If $g(d)$ decays exponentially fast to 0, say $g(d) \asymp a^d$ as $d \to \infty$ for some $a \in (0,1)$, then we get the curse of dimensionality.
\end{itemize}

\item If, for some $c \in (0,1)$, $g(d) \to c$ as $d \to \infty$, then the information complexity does not increase as $d \to \infty$ and may even decrease, i.e., we have strong polynomial tractability and may even have the blessing of dimensionality in this case (our current bounds are not strong enough to show which scenario applies). 

In particular, an interesting choice to consider here would be $\alpha_d = \eta \, C_d$ for some fixed $\eta \in (\sqrt{2}, \sqrt{2\pi}]$. In this case we have $\alpha_d \in (C_d I_d, 1]$ by Lemma~\ref{lem:C.d.I.d.inequalities} and therefore $K_{\eta \, C_d, d}(\bm{x}, \bm{y}) = 1 - \eta^{-1} \|\bm{x} - \bm{y}\|$ is a reproducing kernel for the Sobolev space $H^{(d+1)/2}$.

\end{itemize}




\section{The hemisphere discrepancy and geodesic distance}\label{sec_hemisphere}

In \cite{BDM18} the authors consider a spherical cap discrepancy with respect to hemispheres, which is related to the geodesic distance. This corresponds to the case $\alpha_d \to 0^+$ in the previous section, i.e.
\begin{align*}
K_{\alpha_d, d}(\bm{x}, \bm{y}) & = \int_{-1}^1 \int_{\mathbb{S}^d} \bm{1}_{C(\bm{z}; \alpha_d \, t)}(\bm{x}) \bm{1}_{C(\bm{z}; \alpha_d \, t)}(\bm{y}) \,\mathrm{d} \sigma_d(\bm{z}) \,\mathrm{d} t \\ & \xrightarrow[\alpha_d\to0^+]{}  2 \int_{\mathbb{S}^d} \bm{1}_{C(\bm{z};0)}(\bm{x}) \bm{1}_{C(\bm{z}; 0)}(\bm{y}) \,\mathrm{d} \sigma_d(\bm{z}) = K_{0,d}(\bm{x}, \bm{y}).
\end{align*}


Formula \eqref{Kalpha} is no longer defined in the limit $\alpha_d \to 0^+$. Instead, the corresponding reproducing kernel is given by (using results from \cite{BDM18})
\begin{equation*}
K_{0,d}(\bm{x}, \bm{y}) = 2 \int_{\mathbb{S}^d} 1_{C(\bm{z}; 0)}(\bm{x}) 1_{C(\bm{z}; 0)}(\bm{y}) =   1 - d(\bm{x}, \bm{y}),
\end{equation*}
where $d(\bm{x}, \bm{y})$ is the normalized geodesic distance between $\bm{x}$ and $\bm{y}$. 




The corresponding Stolarsky invariance principle is (see \cite[Theorem~3.1]{BDM18})
\begin{align*}
& \frac{1}{N^2} \sum_{n,m=1}^N K_{0,d}(\bm{x}_n, \bm{x}_m) - \int_{\mathbb{S}^d} \int_{\mathbb{S}^d} K_{0,d}(\bm{x}, \bm{y}) \,\mathrm{d} \sigma_d(\bm{x}) \,\mathrm{d} \sigma_d(\bm{y}) \\
& =  \int_{\mathbb{S}^d} \int_{\mathbb{S}^d} d(\bm{x}, \bm{y}) \,\mathrm{d} \sigma_d(\bm{x}) \,\mathrm{d} \sigma_d(\bm{y}) - \frac{1}{N^2} \sum_{n,m=1}^N d(\bm{x}_n, \bm{x}_m)  \\ & = 2 \int_{\mathbb{S}^d} \left| \frac{1}{2} - \frac{1}{N} \sum_{n=1}^N 1_{C(\bm{x}; 0)}(\bm{x}_n) \right|^2 \,\mathrm{d} \sigma_d(\bm{x})\\
& = (L_2^{\cp,0}(P_{N,d}))^2.
\end{align*}
Note that $\sigma_d(C(\bm{x}; 0)) = 1/2$ since $C(\bm{x}; 0)$ is a hemisphere. In particular, the measure of the hemisphere is independent of the dimension.

We know from \cite{BDM18} that 
\begin{equation*}
\int_{\mathbb{S}^d} \int_{\mathbb{S}^d} d(\bm{x}, \bm{y}) \,\mathrm{d} \sigma_d(\bm{x}) \,\mathrm{d} \sigma_d(\bm{y}) = \frac{1}{2}.
\end{equation*}
Hence, the initial hemisphere discrepancy is $L_2^{\cp,0}(\emptyset) = \frac{1}{\sqrt{2}}$. Using \cite[Theorem~3.2]{BDM18} we obtain that for even $N$, any centrally symmetric point set $P_{N,d}$ satisfies $L_2^{\cp,0}(P_{N,d}) = 0$. This shows that $N^{\cp,0}(\varepsilon, d)= 2$ for any $\varepsilon \in (0,1)$ and $d \in \mathbb{N}$. 

\begin{comment}
\section{The spherical cap $L_p$ discrepancy}

We consider now the $L_p$ norm of the spherical cap discrepancy, which for $p \in [1, \infty)$, is given by
\begin{equation*}
L_{p, \alpha_d}(P_{N,d})  = \left( \int_{-1}^1 \int_{\mathbb{S}^d} \left| \frac{1}{N} \sum_{n=1}^N \bm{1}_{C(\bm{x}, \alpha_d \, t)}(\bm{x}_n) - \sigma_d(C(\bm{x}, \alpha_d \, t)) \right|^p \,\mathrm{d} \sigma_d(\bm{x}) \,\mathrm{d} t \right)^{1/p},
\end{equation*}
and for $p = \infty$ is given by
\begin{equation*}
L_{\infty, \alpha_d}(P_{N,d})  = \sup_{-1 \le t \le 1, \bm{x} \in \mathbb{S}^d} \left| \frac{1}{N} \sum_{n=1}^N \bm{1}_{C(\bm{x}, \alpha_d \, t)}(\bm{x}_n) - \sigma_d(C(\bm{x}, \alpha_d \, t)) \right|.
\end{equation*}
For the initial discrepancy, we have
\begin{equation*}
\frac{1}{2} \le L_{\infty, \alpha_d}(\emptyset) = \sigma_d(C(\bm{x}, \alpha_d)) \le 1.
\end{equation*}
Using the Riesz-Thorin interpolation theorem we obtain for $2 \le p < \infty$ that
\begin{equation*}
L_{p, \alpha_d}(\emptyset) \le \left(1 - \frac{C_d I_d}{\alpha_d} \right)^{1/p}.
\end{equation*}


From the $L_2$ spherical cap discrepancy lower bound \eqref{L2_lower_bound} we have for $2 \le p \le \infty$
\begin{equation*}
L_{p, \alpha_d}(P_{N,d}) \ge \sqrt{C \frac{\sqrt{2} C_d}{\alpha_d} \frac{1}{\sqrt{d} N^{1 + 1/d}}.}
\end{equation*}
Thus
\begin{equation*}
N^{\cp}(\varepsilon, d, p) \gtrsim \left( \left( \frac{\alpha_d}{\alpha_d-C_d I_d} \right)^{2/p} \frac{\sqrt{2} C_d}{\alpha_d} \frac{1}{\sqrt{d} \varepsilon^2} \right)^{\tfrac{d}{d+1}}.
\end{equation*}
For $\alpha_d = \sqrt{2} C_d$ we obtain
\begin{equation*}
N^{\cp}(\varepsilon, d, p) \gtrsim \left( \frac{d^{\tfrac{2}{p} - \tfrac{1}{2}}}{\varepsilon^2} \right)^{\tfrac{d}{d+1}}.
\end{equation*}
This implies that for $2 \le p \le 4$ we do not have a blessing of dimensionality by choosing $\alpha_d = \sqrt{2} C_d$.


\todo{Estimate $\mathbb{E}( L_{p,\alpha_d}^p(P_{N,d})^p)$ for $2 \le p \le \infty$ so that we can get an upper bound for $N^{\cp}(\varepsilon, d, p)$?}
\end{comment}






%\appendix

%\section*{Appendix}

\begin{appendix}
%\section{The distance integral $I_d$}\label{sec_dist_int}
\section{Proof of Lemma~\ref{lem:C.d.I.d.inequalities} and of Eq.~\eqref{asym:Cxt}}\label{sec_dist_int} 

\begin{proof}[Proof of Lemma~\ref{lem:C.d.I.d.inequalities}] 

Both, $I_d$ in \eqref{def:Id} and $C_d$ in \eqref{def:Cd}, can be understood as a function of the real variable $d \in \mathbb{N}$.  
%
The log-convex property of the Gamma function, i.e., $\Gamma(\tfrac{1}{2}(u+v)) < \sqrt{\Gamma(u) \Gamma(v)}$ for $u,v \in \mathbb{R}^+$, will be essential for our proofs of the lower and upper bounds of $C_d$ and $I_d$. 

\paragraph{Lower bound for $C_d$:} 
Log-convexity of the Gamma function gives
\begin{equation}
\label{eq:Gamma.d.2.bound}
\Gamma(\tfrac{d}{2}) 
= \Gamma(\tfrac{1}{2}( \tfrac{d-1}{2} + (\tfrac{d-1}{2} + 1))) 
< \sqrt{\Gamma(\tfrac{d-1}{2})} \, \sqrt{\Gamma(\tfrac{d-1}{2} + 1 )} = \sqrt{\tfrac{d-1}{2}} \, \Gamma(\tfrac{d-1}{2}).
\end{equation}
Hence
\begin{equation*}
C_d = \frac{1}{d} \, \frac{\Gamma( \tfrac{d+1}{2} )}{\sqrt{\pi} \, \Gamma( \tfrac{d}{2} )} > \frac{1}{d} \, \frac{\tfrac{d-1}{2} \ \Gamma(\tfrac{d-1}{2})}{\sqrt{\pi} \, \sqrt{\tfrac{d-1}{2}} \, \Gamma( \tfrac{d-1}{2} )}
\end{equation*}
and the lower bound for $C_d$ follows. 

\paragraph{Upper bound for $C_d$:} Using \eqref{eq:Gamma.d.2.bound} with $d$ replaced by $d+1$ we obtain 
\begin{equation*}
C_d = \frac{1}{d} \frac{\Gamma(\tfrac{d+1}{2})}{\sqrt{\pi} \, \Gamma(\tfrac{d}{2})} < \frac{1}{d} \, \frac{\sqrt{d/2} \, \Gamma(\tfrac{d}{2})}{\sqrt{\pi} \, \Gamma(\tfrac{d}{2})} = \frac{1}{\sqrt{2\pi d}}.
\end{equation*}


\paragraph{Lower bound for $I_d$:} 
We estimate the denominator of \eqref{eq:I.d.value} from above using the log-convexity inequality and the duplication formula \cite[5.5.5]{NIST} for the Gamma function 
\begin{equation*}
\begin{split}
\Gamma( d + \tfrac{1}{2} ) 
&= \Gamma(\tfrac{1}{2}( d + ( d + 1 ) ) ) < \sqrt{\Gamma( d ) } \, \sqrt{\Gamma( d + 1 )} = \sqrt{d} \, \Gamma( d ) \\
&= \sqrt{d} \, 2^{d-1} \, \frac{1}{\sqrt{\pi}} \, \Gamma(\tfrac{d}{2}) \, \Gamma(\tfrac{d+1}{2})
\end{split}
\end{equation*}
and get
\begin{equation*}
I_d 
= 
2^d \, \frac{\Gamma( \tfrac{d + 1}{2} ) \, \Gamma( \tfrac{d + 1}{2} )}{\sqrt{\pi} \, \Gamma( d + \tfrac{1}{2} )}
> 2 \, \frac{\Gamma( \tfrac{d + 1}{2})}{\sqrt{d} \, \Gamma( \tfrac{d}{2} )}. 
\end{equation*}
Applying \eqref{eq:Gamma.d.2.bound}, we arrive at
\begin{equation*}
I_d > 2 \, \frac{\tfrac{d - 1}{2} \ \Gamma( \tfrac{d - 1}{2} )}{\sqrt{d} \, \sqrt{\tfrac{d-1}{2}} \, \Gamma( \tfrac{d-1}{2} )}
\end{equation*}
from which the lower bound follows. 
%\todo{
%The function
%\begin{equation*}
%f( d ) := \frac{I_d}{\sqrt{2} \, \sqrt{1 - \frac{1}{d}}}
%\end{equation*}
%is strictly monotonically decreasing for $d \geq 1$ with 
%\begin{equation*}
%\lim_{d \to \infty} f( d ) = 1. 
%\end{equation*}
%This shows the lower bound of $I_d$ for all real $d\geq 1$.
%}

\paragraph{Upper bound for $I_d$:} The duplication formula for the Gamma function applied to the denominator of $I_d$ gives
\begin{equation} \label{eq:Id.b}
I_d = \sqrt{2} \, \frac{\Gamma(\tfrac{d}{2} + \tfrac{1}{2} ) \, \Gamma( \tfrac{d}{2} + \tfrac{1}{2} )}{\Gamma(\tfrac{d}{2} + \tfrac{1}{4} ) \, \Gamma( \tfrac{d}{2} + \tfrac{3}{4} )}.
\end{equation}
We shall show that the function
\begin{equation*}
f( \alpha ) := \Gamma( \tfrac{d}{2} + \alpha ) \, \Gamma( \tfrac{d}{2} + 1 - \alpha )
\end{equation*} 
has a unique minimum at $\alpha = 1/2$ in $[0,1]$ for each $d \in \mathbb{N}$. This shows the upper bound $I_d < \sqrt{2}$. The first derivative of $f$ can be expressed in terms of the digamma function $\psi( s ) := \Gamma^\prime(s) / \Gamma( s )$ for $s \neq 0, -1, -2, \dots$, i.e.
\begin{equation*}
\begin{split}
f^\prime( \alpha ) 
&= f( \alpha ) \left( \psi(\tfrac{d}{2} + \alpha ) - \psi( \tfrac{d}{2} + 1 - \alpha \right) \\
&= f( \alpha ) \, \sum_{n=0}^\infty \frac{2\alpha-1}{(n + \tfrac{d}{2} + 1 - \alpha ) \, ( n + \tfrac{d}{2} + \alpha )}, 
\end{split}
\end{equation*}
where we used the series expansion \cite[5.7.6]{NIST}
\begin{equation*}
\psi( s ) = - \gamma + \sum_{n=0}^\infty \left( \frac{1}{n+1} - \frac{1}{n + s} \right) \qquad \mbox{for } s \neq 0, -1, -2, \dots
\end{equation*}
and $\gamma$ denotes the Euler–Mascheroni constant
\begin{equation*}
\gamma = \lim_{n\to \infty} \left( \sum_{k=1}^n \frac{1}{k} - \log n \right).
\end{equation*}
Clearly, $f^\prime( 1/2 ) = 0$ and $f$ is strictly monotonically decreasing on $[0,1/2]$ and strictly monotonically increasing on $[1/2,1]$. Hence, $f(\alpha)$ has a unique minimum at $\alpha = 1/2$ in $[0,1]$. 

Asymptotic expansions for Gamma function ratios (cf. \cite[5.11.13, 5.11.15]{NIST}) applied to \eqref{eq:Id.b} yields
\begin{equation*}
I_d = \sqrt{2} \, \left( \frac{d}{2} \right)^{-\tfrac{1}{4}} \left( 1 - \frac{1}{16 d} + \mathcal{O}\Big( \frac{1}{d^2} \Big) \right) \left( \frac{d}{2} \right)^{\tfrac{1}{4}} \left( 1 - \frac{1}{16 d} + \mathcal{O}\Big( \frac{1}{d^2} \Big) \right) \qquad \text{as $d \to \infty$.}
\end{equation*}
The result follows. 

\paragraph{Lower and upper bound for $C_d \, I_d$:} These bounds follow from multiplying the corresponding bounds of $C_d$ and $I_d$. 

\paragraph{Lower and upper bound for $C_d \, I_d / ( 1 - C_d \, I_d )$:} These bounds follow from the bounds of $C_d \, I_d$ and simplification.
\end{proof}

  
\iffalse
Let $\bm{e}_1 = (1, 0, \ldots, 0)^\top$. Then
\begin{equation*}
I_d = \int_{\mathbb{S}^d} \|\bm{e}_1 - \bm{x} \| \,\mathrm{d} \sigma_d(\bm{x}) = \int_{\mathbb{S}^d} \sqrt{2 ( 1- \langle \bm{e}_1, \bm{x} \rangle ) } \,\mathrm{d} \sigma_d(\bm{x}),
\end{equation*}
where we have used the formula $\|\bm{x} - \bm{y}\| = \sqrt{2 ( 1- \langle \bm{x}, \bm{y} \rangle ) }$ for $\bm{x},\bm{y} \in \mathbb{S}^d$. From \cite[Appendix~A]{BD13a} we have
\begin{equation*}
I_d = \frac{\omega_{d-1}}{\omega_d} \int_{-1}^1 \sqrt{2 (1-t)} (1-t^2)^{d/2-1} \,\mathrm{d} t,
\end{equation*}
where $\omega_d$ denotes the surface area of $\mathbb{S}^d$. Now, from the Beta integral we obtain
\begin{equation*}
B(p, q) = 2^{1-p-q} \int_{-1}^1 (1-t)^{p-1} (1+t)^{q-1} \,\mathrm{d} t.
\end{equation*}
By setting $p = (d+1)/2$ and $q = d/2$ we get
\begin{equation*}
I_d =  2^{d} \frac{\omega_{d-1}}{\omega_d}  B\left(\frac{d+1}{2},\frac{d}{2}\right).
\end{equation*}
Rewriting this in terms of the $\Gamma$ function we obtain
\begin{equation*}
I_d = 2^{d} \frac{\Gamma((d+1)/2)}{\sqrt{\pi} \Gamma(d/2)} \frac{\Gamma((d+1)/2) \Gamma(d/2)}{\Gamma(d+1/2)} = \frac{2^{d}}{\sqrt{\pi}} \frac{ [ \Gamma( (d+1)/2)]^2}{\Gamma(d+1/2)}.
\end{equation*}

Using Stirling's formula $\Gamma(x+1) \sim \sqrt{2\pi x} (x/\mathrm{e})^x$ we obtain that
\begin{equation*}
I_d \to \sqrt{2} \mbox{ as } d \to \infty.
\end{equation*}
More precise estimations (for instance, using the $\log$-Gamma expansion) yield
\begin{equation*}
I_d = \sqrt{2} \left(1 - \frac{1}{8d} + \mathcal{O}\left(\frac{1}{d^2}\right) \right).
\end{equation*}
\fi


\iffalse
Next we show that for every $d\in \mathbb{N}$ we have 
\begin{equation}\label{estId}
I_d < \sqrt{2}.
\end{equation}

For the proof of this estimate we use the following lemma:

\begin{lemma}\label{lem:coords}
For independent random points $X,Y$ distributed according to $\sigma_d$ on $\mathbb{S}^d$ we have $\mathbb{E}[\langle X,Y\rangle] = 0$ and $\mathbb{E}[\langle X,Y \rangle ^2] =\tfrac{1}{d+1}$.
\end{lemma}

\begin{proof}
Let $X=(X_1,\ldots,X_{d+1})$ be a random point distributed according to $\sigma_d$ on $\mathbb{S}^d$. Then for $i\in\{1,\dots,d+1\}$, by odd symmetry (or rotational invariance), $\mathbb{E}[X_i]=0$ and, moreover,
\[
1 = \mathbb{E}[\|X\|^2] = \mathbb{E}\left[\sum_{i=1}^{d+1} X_i^2\right]
= \sum_{i=1}^{d+1} \mathbb{E}[X_i^2].
\]
By symmetry all $\mathbb{E}[X_i^2]$ are equal, hence each equals $1/(d+1)$.

Now, using independence of $X$ and $Y$, we obtain 
\[
\mathbb{E}[\langle X , Y\rangle] = \sum_{i=1}^{d+1} \mathbb{E}[X_i]\,\mathbb{E}[Y_i] = 0
\]
and 
\begin{align*}
\mathbb{E}\bigl[\langle X,Y\rangle^2\bigr]
= & \mathbb{E}\left[\sum_{i,j=1}^{d+1} X_i Y_i X_j Y_j\right]
= \sum_{i,j=1}^{d+1} \mathbb{E}[X_i X_j]\, \mathbb{E}[Y_i Y_j]\\
= & \sum_{i=1}^{d+1} \mathbb{E}[X_i^2]\, \mathbb{E}[Y_i^2]
= (d+1)\left(\frac{1}{d+1}\right)^2
= \frac{1}{d+1},
\end{align*}
where we also used the fact that $\mathbb{E}[X_i X_j]=0$ for $i\neq j$ by symmetry.
\end{proof}

\begin{proof}[Proof of \eqref{estId}]
Let $X,Y$ be independent random points distributed according to $\sigma_d$ on $\mathbb{S}^d$. Then, using concavity of the square-root function and Jensen's inequality
\[
I_d = \mathbb{E}[\|X-Y\|]= \mathbb{E}[\sqrt{2(1-\langle X,Y\rangle)}] \le \sqrt{\mathbb{E}[2(1-\langle X,Y\rangle)]} = \sqrt{2(1-\mathbb{E}[\langle X,Y\rangle])}= \sqrt{2},
\]
according to Lemma~\ref{lem:coords}. Thus $I_d \le \sqrt{2}$.

To show the inequality is strict, recall that equality in Jensen's inequality for a strictly concave function occurs if and only if the argument is almost surely constant. Here that would require $\langle X,Y\rangle$ to be almost surely constant. However, according to Lemma~\ref{lem:coords}, 
\begin{equation*}%\label{eq:varT}
\mathrm{Var}[\langle X,Y\rangle] = \mathbb{E}[\langle X,Y\rangle^2] - (\mathbb{E}[\langle X,Y\rangle])^2 = \frac{1}{d+1} > 0,
\end{equation*}
so $\langle X,Y\rangle$ is not almost surely constant, and therefore $I_d < \sqrt{2}$.
\end{proof}

\begin{remark}\rm 
It follows from \eqref{estId} that $1 - \tfrac{I_d}{\sqrt{2}} > 0$ which ensures that $K_{\alpha,d}$ is a reproducing kernel.
\end{remark}

\fi 

\iffalse
\begin{remark}[Asymptotics]
In high dimensions, two independent uniform points on $\mathbb{S}^d$ are almost orthogonal, so $\langle X,Y\rangle$ concentrates near $0$ and $\|X-Y\|$ concentrates near $\sqrt{2}$. One can show $I_d \uparrow \sqrt{2}$ as $d\to\infty$, while the strict inequality $I_d<\sqrt{2}$ holds for every finite $d$ by the argument above.
\end{remark}
\fi



Finally, we present the proof of \eqref{asym:Cxt}:
\begin{proof}[Proof of \eqref{asym:Cxt}]
Let $t \in (0,1)$ be fixed. By the Funk-Hecke formula \cite{M06} we get
\begin{equation*}
\sigma_d(C(\bm{x};t)) = \frac{\omega_{d-1}}{\omega_d} \int_t^1 (1-\tau^2)^{\frac{d}{2}-1} \mathrm{d} \tau.   
\end{equation*}
Here, $\omega_d$ denotes the surface area of the sphere $\mathbb{S}^d$. In particular
\begin{equation*}
\frac{\omega_{d-1}}{\omega_d} = \frac{\Gamma(\tfrac{d+1}{2})}{\sqrt{\pi}\  \Gamma(\frac{d}{2})} = d \, C_d. 
\end{equation*}
The change of variable $(1-t^2) \, u = 1 - \tau^2$, $0 \leq \tau \leq 1$, with $(1-t^2) \, \mathrm{d} u = - 2 \tau \, \mathrm{d} \tau$ yields
\begin{equation*}
\sigma_d(C(\bm{x};t)) = \frac{d}{2} \, C_d \left( 1 - t^2 \right)^{\frac{d}{2}} \int_0^1 u^{\frac{d}{2}-1} \left( 1 - ( 1 - t^2 ) \, u \right)^{-\frac{1}{2}} \mathrm{d} u. 
\end{equation*}
Integration by parts gives
\begin{equation*}
\sigma_d(C(\bm{x};t)) = C_d \, \frac{( 1 - t^2 )^{\frac{d}{2}}}{t} - \frac{1}{2} \, C_d \left( 1 - t^2 \right)^{\frac{d}{2}} \left( 1 - t^2 \right) \int_0^1 u^{\frac{d}{2}} \left( 1 - ( 1 - t^2 ) \, u \right)^{-\frac{3}{2}} \mathrm{d} u.
\end{equation*}
Since
\begin{equation*}
0 \leq \int_0^1 u^{\frac{d}{2}} \left( 1 - ( 1 - t^2 ) \, u \right)^{-\frac{3}{2}} \mathrm{d} u \leq \frac{1}{t^3} \int_0^1 u^{\frac{d}{2}} \, \mathrm{d} u = \frac{1}{\tfrac{d}{2}+1} \, \frac{1}{t^3} 
\end{equation*}
and using Lemma~\ref{lem:C.d.I.d.inequalities}, we arrive at the two-sided estimates
\begin{equation}
\label{eq:sigma.d.estimate}
\frac{\sqrt{1 - \tfrac{1}{d}}}{\sqrt{2\pi d}} \, \frac{( 1 - t^2 )^{\frac{d}{2}}}{t} \left( 1  - \frac{1}{d+2} \, \frac{1}{t^2} \right) 
<
\sigma_d(C(\bm{x};t))
<
\frac{1}{\sqrt{2\pi d}} \, \frac{( 1 - t^2 )^{\frac{d}{2}}}{t}
\end{equation}
or
\begin{equation*}
\sigma_d(C(\bm{x};t)) = C_d \, \frac{( 1 - t^2 )^{\frac{d}{2}}}{t} \left( 1 + \mathcal{O}\left( \frac{1}{d+2} \, \frac{1}{t^2} \right) \right) \qquad \text{as $d \to \infty$,}
\end{equation*}
where the hidden constant is, indeed, $1$. The result~\eqref{asym:Cxt} follows from the observation
\begin{equation*}
\left( 1 - t^2 \right)^{\frac{d}{2}} 
= 
\exp\Big( \frac{d}{2} \log( 1 - t^2 ) \Big) 
=
\exp\left( - \frac{d}{2} \left( t^2 + \frac{t^4}{2} + \cdots \right) \right)
\end{equation*}
when using the Taylor series expansion of $\log( 1 - u )$ for $u = t^2$ near $0$.
\end{proof}


\section{Proof of Eq.~\eqref{L2_lower_bound}}
\label{app:lower.bound}

\begin{proof}[Proof of Eq.~\eqref{L2_lower_bound}]
We apply the method in \cite{BB25}. According to \cite[Eq.~(5)]{BB25} we have the expansion 
\begin{equation*}
\frac{1}{\sqrt{2}} \| \bm{x} - \bm{y} \| 
=
1 + \sum_{\ell=1}^\infty (-1)^{\ell} \binom{\tfrac{1}{2}}{\ell} \langle \bm{x}, \bm{y} \rangle^\ell,
\end{equation*}
where the coefficients are negative for $\ell\in \mathbb{N}$ and can be given in terms of Pochhammer symbols ($(a)_0 := 1$ and $(a)_{n+1} = (a)_{n} \, (n+ a)$ for $n \in \mathbb{N}_0$); i.e.
\begin{equation*}
(-1)^{\ell} \binom{\tfrac{1}{2}}{\ell} = \frac{(-\frac{1}{2})_{\ell}}{\ell!} = \frac{1}{\Gamma( - \frac{1}{2} )} \, \frac{\Gamma( \ell - \frac{1}{2} )}{\Gamma( \ell + 1 )}. %\sim -\frac{1}{2 \sqrt{\pi }} \, \fritz{\frac{1}{\ell^{3/2}}} \qquad \text{as $\ell \to \infty$.}
\end{equation*}
Substitution into Stolarsky's invariance principle \eqref{eq:Stolarsky.gen} taking into account \eqref{Kalpha} gives
\begin{eqnarray*}
\lefteqn{(L_2^{\cp,\alpha_d}(P_{N,d}))^2}  \\
& = & \frac{1}{N^2} \sum_{n,m=1}^N K_{\alpha_d,d}(\bm{x}_n, \bm{x}_m) - \int_{\mathbb{S}^d} \int_{\mathbb{S}^d} K_{\alpha_d,d}(\bm{x}, \bm{y}) \,\mathrm{d} \sigma_d(\bm{x}) \,\mathrm{d} \sigma_d(\bm{y}) \\
& = &
\frac{\sqrt{2} C_d}{\alpha_d} \sum_{\ell=1}^\infty \frac{-(-\frac{1}{2})_{\ell}}{\ell!} \left( \frac{1}{N^2} \sum_{m,n=1}^N \langle \bm{x}_m, \bm{x}_n \rangle^\ell - \int_{\mathbb{S}^d} \int_{\mathbb{S}^d} \langle \bm{x}, \bm{y} \rangle^\ell \,\mathrm{d} \sigma_d(\bm{x}) \,\mathrm{d} \sigma_d(\bm{y}) \right).
\end{eqnarray*}
The parenthetical expression are all non-negative by \cite[Lemma~1]{BB25}. Thus the series expansions adds up non-negative terms. By removing all odd terms and even terms with $\ell < 2M$, $M \in \mathbb{N}$ to be chosen later, we get 
\begin{eqnarray*}
\lefteqn{(L_2^{\cp,\alpha_d}(P_{N,d}))^2}\\
& \geq &
\frac{\sqrt{2} C_d}{\alpha_d} \sum_{r=M}^\infty \frac{-(-\frac{1}{2})_{2r}}{(2r)!} \left( \frac{1}{N^2} \sum_{m,n=1}^N \langle \bm{x}_m, \bm{x}_n \rangle^{2r} - \int_{\mathbb{S}^d} \int_{\mathbb{S}^d} \langle \bm{x}, \bm{y} \rangle^{2r} \,\mathrm{d} \sigma_d(\bm{x}) \,\mathrm{d} \sigma_d(\bm{y}) \right).
\end{eqnarray*}
By \cite[Lemma~2]{BB25} 
\begin{equation*}
\int_{\mathbb{S}^d} \int_{\mathbb{S}^d} \langle \bm{x}, \bm{y} \rangle^{2r} \,\mathrm{d} \sigma_d(\bm{x}) \,\mathrm{d} \sigma_d(\bm{y})  
=  
\frac{(\frac{1}{2})_{r}}{(\frac{d+1}{2})_{r}} 
=
\frac{\Gamma(\tfrac{d+1}{2})}{\Gamma(\tfrac{1}{2})} \, \frac{\Gamma(r+\tfrac{1}{2})}{\Gamma(r+\tfrac{d+1}{2})}
\sim
\frac{\Gamma( \frac{d+1}{2} )}{\Gamma( \frac{1}{2} )} \, \frac{1}{r^{\frac{d}{2}}} \quad \text{as $r \to \infty$.}    
\end{equation*}
Furthermore, the integrals are strictly monotonically decreasing in $r$ and the asymptotic term is a strict upper bound for all $r\in \mathbb{N}$. Using this and keeping only the diagonal terms in the double-sum we obtain  
\begin{equation*}
\frac{1}{N^2} \sum_{m,n=1}^N \langle \bm{x}_m, \bm{x}_n \rangle^{2r} - \int_{\mathbb{S}^d} \int_{\mathbb{S}^d} \langle \bm{x}, \bm{y} \rangle^{2r} \,\mathrm{d} \sigma_d(\bm{x}) \,\mathrm{d} \sigma_d(\bm{y}) 
\geq 
\frac{1}{N} - \frac{(\frac{1}{2})_{M}}{(\frac{d+1}{2})_{M}} 
\geq 
\frac{1}{N} - \frac{\Gamma( \frac{d+1}{2} )}{\Gamma( \frac{1}{2} )} \, \frac{1}{M^{d/2}}.
\end{equation*}
By \cite[Eq.~(25)]{BB25} we have
\begin{equation*}
\sum_{r=M}^\infty \frac{-(-\frac{1}{2})_{2r}}{(2r)!} 
\geq 
\frac{1}{2 \sqrt{2 \pi}} \, \frac{1}{M^{1/2}} \qquad \mbox{for all } M \in \mathbb{N},
\end{equation*}
where the right-hand side is indeed the leading term of the large $M$-asymptotics. Thus we get 
\begin{equation*}
(L_2^{\cp,\alpha_d}(P_{N,d}))^2  \geq 
\frac{\sqrt{2} C_d}{\alpha_d} \frac{1}{2 \sqrt{2 \pi}}  \left( \frac{1}{N} - \frac{\Gamma( \frac{d+1}{2} )}{\Gamma( \frac{1}{2} )} \, \frac{1}{M^{d/2}} \right) \frac{1}{M^{1/2}} = \frac{\sqrt{2} C_d}{\alpha_d} \frac{1}{2 \sqrt{2 \pi}} \, \frac{F( \frac{M^{d/2}}{N} )}{N^{1+\frac{1}{d}}} 
\end{equation*}
provided $M$ is such that the parenthetical expression is positive. The function 
\begin{equation*}
F(x) := \frac{x - \beta}{x^{1+\gamma}}, \qquad x \geq \beta := \frac{\Gamma( \frac{d+1}{2} )}{\Gamma( \frac{1}{2} )}, \qquad \gamma := \frac{1}{d}
\end{equation*}
has a unique maximum at $x^* = \beta ( 1 + \frac{1}{\gamma} ) > \beta$ with value 
\begin{equation} 
\label{eq:F.x.star}
F(x^*) = \frac{1}{( 1 + \gamma ) \big( \beta ( 1 + \frac{1}{\gamma} ) \big)^{\gamma}}=\frac{d}{(d+1)^{1+\frac{1}{d}}} \left(\frac{\Gamma(\frac{1}{2})}{\Gamma(\frac{d+1}{2})}\right)^{\frac{1}{d}}.
\end{equation}
This, together with the fact that $\sqrt{2}> I_d$, proves \eqref{L2_lower_bound}.
\end{proof}

\end{appendix}

\newpage

\begin{thebibliography}{99}

\bibitem{BB25} D. Bilyk, J. S. Brauchart: On the lower bounds for the spherical cap discrepancy. \url{arXiv:2502.15984}

\bibitem{BD} D. Bilyk, F. Dai: Geodesic distance Riesz energy on the sphere. Trans. Amer. Math. Soc. 372(5): 3141--3166, 2019.

\bibitem{BDM18} D. Bilyk, F. Dai, R. Matzke: The Stolarsky principle and energy optimization on the sphere. Constr. Approx. 48(1): 31--60, 2018.

\bibitem{BD13a} J. S. Brauchart, J. Dick: A simple proof of Stolarsky's invariance principle. Proc. Amer. Math. Soc. 141(6): 2085--2096, 2013.

\bibitem{BD13b} 
J. S. Brauchart, J. Dick: A characterization of Sobolev spaces on the sphere and an extension of Stolarsky's invariance principle to arbitrary smoothness. Constr. Approx. 38(3): 397--445, 2013.

\bibitem{BSSW14}
J. S. Brauchart, E. B. Saff, I. H. Sloan, R. S. Womersley: QMC designs: optimal order quasi Monte Carlo integration schemes on the sphere. Math. Comp. 83(290): 2821--2851, 2014.

\bibitem{ChS13} C. Choirat, R. Seri: Numerical properties of generalized discrepancies on spheres of arbitrary dimension, J. Complexity 29(2): 216--235, 2013.

\bibitem{DP14a} J. Dick, F. Pillichshammer: Discrepancy theory and quasi-Monte Carlo integration. In: {\em A Panorama of Discrepancy Theory} (W. Chen, A. Srivastav, G. Travaglini, eds.), pp. 539--619, Lecture Notes in Math., No.~2107, Springer, Cham, 2014.

\bibitem{M06} C. Müller: {\em Spherical Harmonics.} Lecture Notes in Math., No.~17, Springer, Berlin, 1966.

\bibitem{NIST}
NIST Digital Library of Mathematical Functions. \url{https://dlmf.nist.gov/}, Release 1.2.4 of 2025-03-15. F. W. J. Olver, A. B. Olde Daalhuis, D. W. Lozier, B. I. Schneider, R. F. Boisvert, C. W. Clark, B. R. Miller, B. V. Saunders, H. S. Cohl, and M. A. McClain, eds.

\bibitem{NP25a} E. Novak, F. Pillichshammer: Intractability results for integration in tensor product spaces. J. Complexity 86: Paper No. 101901, 28 pp., 2025.

\bibitem{NP25b} E. Novak, F. Pillichshammer: The $L_p$-discrepancy for finite $p>1$ suffers from the curse of dimensionality. Proc. Amer. Math. Soc. 153(4): 1447--1459, 2025.

\bibitem{NW08} E. Novak, H. Wo\'{z}niakowski: {\em Tractability of Multivariate Problems. Volume I: Linear Information.} EMS Tracts in Mathematics 16, Z\"urich, 2008.

\bibitem{NW10} E. Novak, H. Wo\'{z}niakowski: {\em Tractability of Multivariate Problems. Volume II: Standard Information for Functionals.} EMS Tracts in Mathematics 12, Z\"urich, 2010.

\bibitem{Wo99} H. Wo\'{z}niakowski: Efficiency of quasi-Monte Carlo algorithms for high dimensional integrals. In: {\em Monte Carlo and Quasi-Monte Carlo Methods 1998} (H. Niederreiter and J. Spanier, eds.), pp.~114--136, Springer Verlag, Berlin, 2000.

\end{thebibliography}

\end{document}
